%% ============================================================
%% VALIDATION PROTOCOL -- split 2 of 3.
%% Drop-in for <<KEEP:stage1>> <<KEEP:stage2>>.
%% Insert BEFORE the existing drafted "validation resolution bounds claimable fidelity"
%% paragraph, which this file does not duplicate and which follows on naturally.
%% Source: v1 lines 1138-1201.
%%
%% NOTE: the gate values quoted here must match tab_val_targets. Per
%% VALIDATION_APPROACH.md the ex-ante gates should be restored to their v1 values
%% (return-temperature MAE 2.5 K, not 1.0 K) before this is merged.
%% ============================================================

The network was upgraded with the heat pump, electrode boiler and thermal storage
\emph{after} the measurement record was collected, so validation follows a two-stage
strategy after \citet{Kus2025} and \citet{Maldonado2024}.

\paragraph{Stage 1: direct network validation.}
The pipe model is validated against pre-upgrade monitoring data with the new assets set to
zero capacity, calibrated by the branch-sequential $U_p$-tuning procedure of
\citet{Maldonado2024} but constrained to a defensible coefficient range
for the pipe class. Table~\ref{tab:valtargets} lists the key performance indicators and
their gates, fixed \emph{ex ante} from published practice before the calibration was run.

The use of these gates is deliberate. All are reported in
Section~\ref{subsec:validation}, including the several that are not met. In particular, no
model of this network can meet the temperature gates because of the instrumentation
limitations described there. Retaining gates that the instrumentation cannot support and
reporting the resulting failures is therefore the point rather than an oversight: these
failures support the claim that the resolution of a model cannot be verified beyond that of
its metering. The conclusions are instead anchored to annual delivered energy, for which the
gate is met, and hence to the model-implied annual network loss that it determines, the
quantity on which the cost results depend.

\paragraph{Stage 2: necessary-condition verification.}
Dispatch of the new assets is verified against plausibility bounds rather than against
measurements: the heat-pump coefficient of performance must remain within \([2.5,5.5]\),
the thermal-storage state of charge within \([0.05,0.95]\bar{E}_s\), simultaneous charging
and discharging are prohibited, and per-step energy closure is required. These checks are
necessary but not sufficient conditions. A post-upgrade measurement campaign is therefore
the natural follow-up and is identified as a limitation.

\paragraph{Structural sanity checks.}
Two checks confirm internal consistency. The nonlinear reference with its quadratic
constraints deactivated reproduces the baseline objective to within $0.01$\,\%; and the
relaxation property $\mathrm{cost}(\CP) \le \mathrm{cost}(\Lone)$ holds on every one of the
135 synthetic instances, since a copperplate can never cost more than the resolved model of
the same network. A violation of either would indicate a formulation error rather than a
finding.
